% =========================================================
% Calculus — Notes Index
% =========================================================

% ---------------------------------------------------------
% Breadcrumb
% ---------------------------------------------------------
\begin{tcolorbox}[
  colback=gray!6,
  colframe=gray!40,
  arc=2pt,
  left=8pt, right=8pt, top=6pt, bottom=6pt,
  title={\small\textbf{Where You Are in the Journey}},
  fonttitle=\small\bfseries
]
\begin{center}
\small
Real Numbers
$\;\to\;$ \textbf{Calculus (Computational)}
$\;\to\;$ Real Analysis
$\;\to\;$ Multivariable Analysis
$\;\to\;$ $\cdots$
\end{center}

\medskip
\noindent\textbf{How we got here.}
The real number system is in place, together with the algebraic
and order-theoretic properties that make it work.
Computational calculus takes those numbers as given and builds
the practical machinery of rates of change and accumulation ---
the toolkit that physicists, engineers, and applied
mathematicians use daily.

\medskip
\noindent\textbf{What this chapter builds.}
The three-semester undergraduate calculus sequence: limits and
continuity (Calc~1), differentiation rules and techniques,
the definite integral and the Fundamental Theorem of Calculus,
applications; then sequences and series, convergence tests,
and power series (Calc~2); then vectors, partial derivatives,
multiple integrals, and the theorems of Green, Stokes, and
Gauss (Calc~3).
The emphasis throughout is on computation and technique rather
than on $\varepsilon$-$\delta$ proof.

\medskip
\noindent\textbf{Relationship to Analysis.}
This chapter is a \emph{companion} to the analysis chapters,
not a replacement.
Real Analysis (bounding, sequences, continuity, integration)
provides the rigorous foundations that justify why the
computational rules work.
This chapter provides the fluency with those rules that makes
analysis easier to read and motivates why the foundations
matter.
It is useful to have both in hand: compute freely here,
understand why it works in analysis.

\medskip
\noindent\textbf{Where this leads.}
Multivariable calculus feeds directly into differential geometry
(the derivative as a linear map, the Jacobian, differential
forms) and into vector analysis (div, grad, curl, and the
integral theorems).
Series and power series feed into complex analysis and the
study of analytic functions.
\end{tcolorbox}
\vspace{1em}

% ---------------------------------------------------------
% Structural Role
% ---------------------------------------------------------
\begin{tcolorbox}[
  colback=gray!6,
  colframe=gray!40,
  arc=2pt,
  left=8pt, right=8pt, top=6pt, bottom=6pt,
  title={\small\textbf{Structural Role: Calculus --- Fluency Before Rigour}},
  fonttitle=\small\bfseries
]
Computational calculus occupies an unusual position in this
curriculum: it is both a prerequisite and a companion.

As a prerequisite, the differentiation and integration
techniques developed here are assumed throughout the analysis
chapters.
It would be difficult to appreciate the Fundamental Theorem of
Calculus as a theorem if one has never computed a definite
integral as an antiderivative.

As a companion, the calculus sequence provides a running
source of examples.
Every convergence theorem in analysis has a concrete calculus
instance.
Every abstract definition of derivative or integral in analysis
has a computational counterpart here that makes the definition
feel natural rather than arbitrary.

The guiding principle is: \emph{compute first, prove later}.
This chapter is the ``compute'' half of that programme.
\end{tcolorbox}
\vspace{1em}

% ---------------------------------------------------------
% Structural Roadmap
% ---------------------------------------------------------
\subsection*{Structural Roadmap}

\begin{center}
\textbf{Technique $\longrightarrow$ Application $\longrightarrow$ Extension}
\end{center}

The progression follows the standard three-semester sequence:

\medskip
\noindent\textbf{Calculus I --- Single-Variable Differential Calculus}
\begin{enumerate}
  \item Limits: intuitive definition, limit laws, one-sided limits,
        limits at infinity
  \item Continuity: definition, types of discontinuity, the
        Intermediate Value Theorem (statement and use)
  \item The derivative: definition as a limit, geometric meaning,
        differentiability vs.\ continuity
  \item Differentiation rules: power, product, quotient, chain
  \item Derivatives of standard functions: polynomials,
        exponentials, logarithms, trigonometric and inverse
        trigonometric functions
  \item Implicit differentiation and related rates
  \item Linear approximation and differentials
  \item Applications: curve sketching, optimisation, the Mean
        Value Theorem (statement and use), L'H\^opital's rule
\end{enumerate}

\medskip
\noindent\textbf{Calculus II --- Integration and Series}
\begin{enumerate}
  \setcounter{enumi}{8}
  \item The definite integral: Riemann sums, the Fundamental
        Theorem of Calculus
  \item Antiderivatives and indefinite integrals; basic
        integration rules
  \item Integration techniques: substitution, integration by
        parts, trigonometric integrals and substitution,
        partial fractions
  \item Improper integrals: convergence and divergence
  \item Applications of integration: area, volume (discs,
        shells), arc length, surface area
  \item Sequences: convergence, standard limits
  \item Series: geometric and $p$-series; convergence tests
        (comparison, ratio, root, integral, alternating)
  \item Power series: radius of convergence, Taylor and
        Maclaurin series, common expansions
\end{enumerate}

\medskip
\noindent\textbf{Calculus III --- Multivariable and Vector Calculus}
\begin{enumerate}
  \setcounter{enumi}{16}
  \item Vectors in $\mathbb{R}^2$ and $\mathbb{R}^3$: dot
        product, cross product, lines and planes
  \item Vector-valued functions: curves, arc length, curvature,
        the Frenet--Serret frame
  \item Functions of several variables: limits, continuity,
        partial derivatives, the gradient
  \item Directional derivatives, the chain rule, implicit
        differentiation in several variables
  \item Higher-order partial derivatives; the second-derivative
        test; Lagrange multipliers
  \item Double and triple integrals; change of variables;
        polar, cylindrical, and spherical coordinates
  \item Line integrals: scalar and vector fields, work
  \item Surface integrals: flux
  \item Green's theorem, Stokes' theorem, and the Divergence
        theorem
\end{enumerate}

\begin{remark*}[Primary sources]
Stewart, \textit{Calculus: Early Transcendentals};
Strang, \textit{Calculus} (MIT OpenCourseWare, freely available);
Apostol, \textit{Calculus}, Vols.~1--2 (for a more rigorous
treatment that bridges toward analysis).
\end{remark*}



% Future sections (uncomment as content is written):
%
% --- Calculus I ---
% =========================================================
% Calculus — Limits: Section Orchestrator
% =========================================================

\input{volume-vi/calculus/notes/limits/difference-equation}
\input{volume-vi/calculus/notes/limits/limits}
\input{volume-vi/calculus/notes/limits/limit-algebra}
\input{volume-vi/calculus/notes/limits/infinite-limits}










% \input{volume-vi/calculus/notes/notes-limits}
% =========================================================
% Calculus — Continuity: Section Orchestrator
% =========================================================

\section{Continuity}

\input{volume-vi/calculus/notes/continuity/continuity}

% \input{volume-vi/calculus/notes/notes-derivative-definition}
% \input{volume-vi/calculus/notes/notes-differentiation-rules}
% \input{volume-vi/calculus/notes/notes-standard-derivatives}
% \input{volume-vi/calculus/notes/notes-applications-derivatives}
%
% --- Calculus II ---
% \input{volume-vi/calculus/notes/notes-definite-integral}
% \input{volume-vi/calculus/notes/notes-integration-techniques}
% \input{volume-vi/calculus/notes/notes-improper-integrals}
% \input{volume-vi/calculus/notes/notes-applications-integration}
% \input{volume-vi/calculus/notes/notes-sequences}
% \input{volume-vi/calculus/notes/notes-series}
% \input{volume-vi/calculus/notes/notes-power-series}
%
% --- Calculus III ---
% \input{volume-vi/calculus/notes/notes-vectors}
% \input{volume-vi/calculus/notes/notes-vector-functions}
% \input{volume-vi/calculus/notes/notes-partial-derivatives}
% \input{volume-vi/calculus/notes/notes-multiple-integrals}
% \input{volume-vi/calculus/notes/notes-line-surface-integrals}
% \input{volume-vi/calculus/notes/notes-integral-theorems}
